\documentclass{article}
\usepackage[T1]{fontenc}
\usepackage{lmodern}
\usepackage{graphicx}
\usepackage{amsmath,amssymb}
\usepackage[a4paper,margin=2cm]{geometry}
\usepackage[absolute,overlay]{textpos}
\setlength{\TPHorizModule}{1mm}
\setlength{\TPVertModule}{1mm}
\usepackage[indonesian]{babel}
\usepackage{hyperref}
\setlength{\emergencystretch}{2em}
% unit: Halaman 44

\begin{document}

% === Halaman 44 (python/native) ===

Elliptic Curve Elgamal (ECEG)

• Elliptic Curve Elgamal: sistem kriptografi kurva eliptik yang diadopsi dari 

algoritma El Gamal.

• Misalkan Alice ingin mengirim Bob pesan M yang dienkripsi.

• Baik Alice dan Bob menyepakati kurva elelitik dan titik basis B.

• Alice dan Bob membuat kunci privat/kunci publik.

• Alice

• Kunci privat = a

• Kunci publik = PA = a  B

• Bob

• Kunci privat = b

• Kunci publik = PB = b  B

• Alice mengambil plainteks, M, lalu mengkodekannya menjadi sebuah titik, PM, 

pada  kurva eliptik

44

*) Sumber bahan: Debdeep Mukhopadhyay, Elliptic Curve Cryptography ,

 Dept of Computer Sc and Engg IIT Madras

\end{document}
